\chapter{\textbf{How to, and other questions}}
\label{chap:questions}

This last chapter is a troubleshooting list. It collects the errors that
students actually hit in the first weeks, in roughly the order in which they
hit them. When something goes wrong, look here before rerunning the
computation with different numbers.

\section*{The program will not start}
\addcontentsline{toc}{section}{\texorpdfstring{\hskip14mm}{} The program will not start}

\begin{description}[style=unboxed,leftmargin=0pt]

\item[\texttt{castem: command not found}] The executable carries the version
  number: it is \texttt{castem26}, not \texttt{castem}. Check which version
  you installed, and note that a script written for an earlier release will
  call \texttt{castem22} or \texttt{castem19}.

\item[The program starts and immediately exits] There is probably a
  \op{fin;} somewhere in the file, or the file ends in the middle of a
  statement. Remember that a statement is terminated by \texttt{;}, not by a
  newline, so a missing semicolon makes the reader swallow the following
  lines.

\end{description}

\section*{The file will not read}
\addcontentsline{toc}{section}{\texorpdfstring{\hskip14mm}{} The file will not read}

\begin{description}[style=unboxed,leftmargin=0pt]

\item[Errors on a line that looks perfectly correct] Look for a tab
  character. Tabs, and the invisible characters some editors insert, break
  the reader. Configure your editor to insert spaces (see
  section~\ref{sec:intro:editing}) and to show whitespace.

\item[Everything after a certain line is ignored] A statement may span
  several lines but not more than 500 characters in total. Check that you
  have not lost a \texttt{;}.

\item[\texttt{OPERATEUR INCONNU}] Either a genuine typo, or an object whose
  name shadows an operator. Object names must not have the same first four
  characters as an operator: calling a variable \texttt{sigma} destroys
  \op{sigm[a]} for the rest of the session. Restart the program after such a
  mistake -- the damage persists in the session.

\end{description}

\section*{The computation runs but the result is wrong}
\addcontentsline{toc}{section}{\texorpdfstring{\hskip14mm}{} The result is wrong}

\begin{description}[style=unboxed,leftmargin=0pt]

\item[Results off by a factor that looks arbitrary] Operator precedence.
  \textsc{gibiane} evaluates strictly left to right, so
  \op{a + b * c} is \op{(a + b) * c}. Bracket every algebraic expression.
  This is the most common cause of a result that is wrong but not absurd.

\item[Displacements are enormous, or the solver fails] The structure is not
  properly restrained -- a rigid body motion is left free. Count your
  conditions: in 2D you need to block three degrees of freedom at least.

\item[Stresses look right in the middle and wrong at the boundary] The mesh
  is too coarse near a stress concentration, or you are reading values at
  nodes that were extrapolated from the Gauss points. Compare with
  \op{chan 'CHPO'} values and refine.

\item[Nothing changes when I refine the mesh] Check that the parameter you
  changed is actually used by every meshing command. In an early version of
  this very booklet, three of the seven lines of the plate were meshed with a
  variable that was never defined.

\item[A field operation gives an empty or zero result] Two fields can only be
  combined component by component, by name. \texttt{'UX'} times
  \texttt{'SCAL'} is empty, not a product. Rename with \op{exco} first.

\item[\texttt{trace} complains or draws nothing] A \op{chpoint} is plotted
  with the \emph{mesh}, a \op{mchaml} with the \emph{model}. Getting this
  backwards is the usual cause.

\end{description}

\section*{Memory, size and speed}
\addcontentsline{toc}{section}{\texorpdfstring{\hskip14mm}{} Memory, size and speed}

\begin{description}[style=unboxed,leftmargin=0pt]

\item[\texttt{PAS ASSEZ DE PLACE EN MEMOIRE}] Increase the
  \texttt{ESOPE\_PARAM} environment variable, but stay below physical RAM.
  See section~\ref{sec:intro:install}.

\item[The run is very slow] Check that the stiffness matrix is assembled and
  factorised \emph{outside} any loop that does not change it. For large
  three-dimensional models, enable parallelism with \op{opti 'PARA' VRAI;}
  and \op{opti 'ASSI' n;}.

\item[Two runs interfere with each other] \texttt{Cast3M} writes
  \texttt{fort.3} and \texttt{fort.98} in the working directory. Run each
  case in its own directory.

\end{description}

\section*{How to save a standard display}
\addcontentsline{toc}{section}{\texorpdfstring{\hskip14mm}{} How to save a display}

In the graphics window, click \emph{Softcopy} and then one of the
\textsc{postscript} formats; the figures accumulate in \texttt{fort.24} and
are copied to a \texttt{.ps} file when the run ends. To do the same without
touching the window -- which is what you want in a batch run -- use
\op{opti trac 'PS';} and \op{opti ftra 'figure.ps';}, and
\op{opti trac 'X';} to switch back to the screen. See
section~\ref{sec:io:graphics} for the details, including how to change the
grey levels afterwards.

For anything you intend to put in a report, consider exporting to
\textsc{Paraview} with \op{sort 'VTK'} or plotting in Python instead
(chapter~\ref{chap:python}): both give vector output and far better control.

\section*{Where to get help}
\addcontentsline{toc}{section}{\texorpdfstring{\hskip14mm}{} Where to get help}

In order:

\begin{enumerate}
\item \op{info \textit{operator};} in the program, or the notice of that
  operator on the website -- the authoritative description of what the
  arguments mean;
\item the \textbf{Exemples} page of the website: roughly 1\,600 searchable
  example files, and usually the fastest way to see an operator used in
  anger;
\item the \texttt{dgibi} directory of your own installation, which holds the
  test cases;
\item the user mailing list \texttt{cast3m-util@\allowbreak umontpellier.fr},
  where questions get answered;
\item \texttt{support-cast3m@cea.fr} for problems with the code itself.
\end{enumerate}

\noindent The URLs are collected in section~\ref{sec:intro:website}.
